\begin{table}[!h]
\centering
\caption{Gelbach decomposition of the change in the Black pay gap \label{tab:pay-gap-gelbach}}
\centering
\begin{threeparttable}
\begin{tabular}[t]{lrrr}
\toprule
 & Log points & Bootstrap SE & Share of total change\\
\midrule
\addlinespace[0.3em]
\multicolumn{4}{l}{\textbf{Coefficient on P(Black)}}\\
\hspace{1em}P(Black) (draft-free), Table \ref{tab:pay-gap-veteran} column (9): position $\times$ year and prior FE (no draft bucket) & 0.098 & (0.065) & \\
\hspace{1em}P(Black) (draft-free), column-(5) controls & 0.000 & (0.054) & \\
\addlinespace[0.3em]
\multicolumn{4}{l}{\textbf{Contribution of each block to the change}}\\
\hspace{1em}Career stage (A) & 0.051 & (0.021) & \\
\hspace{1em}Prior-season production (B) & 0.038 & (0.034) & \\
\hspace{1em}Career production (C) & 0.003 & (0.027) & \\
\hspace{1em}Pre-NFL signals (D) & 0.004 & (0.033) & \\
\hspace{1em}Total change, $\hat\beta_1^{(1)} - \hat\beta_1^{(5)}$ & 0.097 & (0.063) & \\
\bottomrule
\end{tabular}
\begin{tablenotes}[para]
\item \textit{Notes:} 
\item This table decomposes the change in the coefficient on P(Black) in equation (1) between the raw benchmark (column (9) of Table \ref{tab:pay-gap-veteran}) and the column-(5) controls following \citet{gelbach2016covariates}. Both regressions use the draft-free prediction and its prior's covariates (position at NFL entry, entry era, college type, whether a home county is known; no draft-round bucket), so draft slot and draft round enter only through block D. The column-(5) coefficient on the draft-free prediction (0.000) therefore differs from the main estimate in column (5) of Table \ref{tab:pay-gap-veteran} (-0.004), which uses the primary prediction. The sample is the column-(5) estimation sample of 5,062 freely bargained veteran contracts (UFA and extensions) signed 2014-2026; the outcome is log APY. For each block $k$, the contribution is the coefficient on P(Black) from a regression of the block's fitted contribution $X_k\hat\pi_k$ (coefficients from the column-(5) model) on P(Black) and P(other race) and the fixed effects of the base regression (position $\times$ year signed and the race-prior covariates); the contributions sum to the total change. A positive contribution means the block's controls account for part of a positive raw gap. Shares of the total change are not reported because the total change is not distinguishable from zero (its absolute value is below two bootstrap standard errors). Career stage (A): experience-bin dummies (0, 1, 2, 3, 4-6, 7+ seasons), age at signing and its square, a left-censored-experience indicator and an extension dummy (reference: unrestricted free agents). Whether a team extends a player is itself a team decision. Prior-season production (B): season-before-signing box-score and PFR measures with position-group-specific slopes (passing for QBs, rushing and receiving for RBs, receiving for WRs and TEs, tackles, sacks, QB hits, tackles for loss and pressures for DL and LB, tackles, interceptions, passes defended and coverage allowed for DBs, kicking and punting), games played and injury weeks interacted with position group, and a no-prior-season indicator. Career production (C): career games and injury weeks interacted with position group, and career sums of the block B production measures with the same position-group slopes. Pre-NFL signals (D): log draft pick, an undrafted indicator, draft-round dummies, 247 rating and stars, combine measures (40, vertical, bench, broad jump, cone, shuttle, height, weight) and athletic scores (a composite and size, speed, explosion and agility subscores built from the combine measures) interacted with position group, and the final college team's Power-conference, SRS and HBCU indicators. Controls with incomplete coverage are set to zero with a missing indicator (interacted like the control). Regression calibration identifies the Black-white gap only if, in addition to names and hometown being unrelated to pay given race and the controls, (i) P(Black) is calibrated given every control in the column, not only given the prior's covariates, and (ii) the gap does not vary with the controls (otherwise the coefficient is a variance-weighted average of gaps). Controls that predict race beyond the prior violate (i) and attenuate the coefficient on P(Black) even when they do not affect pay, so part of any change in the coefficient as controls are added is mechanical. The block contributions therefore mix the controls' relation to pay with their information on race; block D (combine, recruit and college signals) is the most likely to carry race information beyond the prior. Within position $\times$ year and prior-covariate cells the column-(5) controls predict P(Black): a joint Wald test of the block A-D controls in a regression of P(Black) on them gives F = 3.44 (p = 0.000), and of the block D controls alone F = 2.64 (p = 0.000) (player-clustered covariance, inverted by pseudo-inverse; numerator degrees of freedom are its rank, 409 of 410 and 259 of 259 coefficients). P(Black) is the predicted probability of being non-Hispanic Black alone; Black Hispanic and multiracial Black players enter P(other race), and Table \ref{tab:pay-gap-race-measures} also reports estimates with P(Black) + P(multiracial) as the Black regressor. The predicted Black share is not calibrated in levels: it lies below the TIDES African-American player share through 2016 and above the self-identified Black-alone share (below the Black-or-multiracial share) from 2019 (Table \ref{tab:race-prediction-tides}); Table \ref{tab:pay-gap-race-measures} reports estimates with the TIDES-raked posterior. Standard errors treat the predicted probabilities as known; the prior is estimated by EM on the whole player population of the database, so the omitted first-stage variance is likely small. Standard errors come from a player-cluster bootstrap with 199 replications: players are resampled with replacement and the whole decomposition is repeated in each replication. P(Black) and P(other race) enter together, so the coefficient on P(Black) compares a Black with a white player. Race is predicted, not observed: each person's probability of being non-Hispanic Black combines first name, surname and hometown county (BIFSG) with an NFL-specific prior estimated by EM on predetermined characteristics (players: position at entry, rookie era, draft round, college type, county availability; staff: role, unit and era at first appearance); documented race statements are not used (notes/race-prediction-design.md). Regressions use the probabilities (regression calibration) and control for the prior's covariates; the coefficient on P(Black) is the Black-white gap if the probabilities are calibrated given the regression's controls and names and hometown are unrelated to the outcome given race and the controls. The validation exhibit compares the predictions with published TIDES shares.
\end{tablenotes}
\end{threeparttable}
\end{table}
